\shorttitle{Min Cut-Path Problem}

\shortauthors{N. Micheľ}

\title[mode = title]{Min Cut-Path Problem}

\author[1]{Norbert Micheľ}

\cormark[1]

\ead{5344553@upjs.sk}

\affiliation[1]{organization={Institute of Computer Science, Faculty of Science, Pavol Jozef Šafárik University in Košice},
            addressline={Jesenná 5},
            postcode={040 01},
            postcodesep={},
            city={Košice},
            country={Slovakia}}

\cortext[1]{Corresponding author}

% The class writes the abstract body verbatim to a file and reads it back. The optional argument (the class's
% own heading) is given explicitly, so the class does not look ahead for one.
\begin{abstract}[\abstractname]
A cut-path between two vertices $u$ and $v$ of a graph is a set of edges that contains both a $u$--$v$ path and a cut separating $u$ from $v$.
It models two servers that communicate along the path while every other route between them uses one of the chosen edges.
We study the \MinCutPath{} problem, which asks for a cut-path with the fewest edges.
We prove that its decision version is NP-complete.
In graphs of diameter two and in graphs in which every two vertices can be separated by at most two edges, a minimum cut-path has $c(u,v) + d(u,v) - 1$ edges.
Here $c(u,v)$ is the size of a minimum cut and $d(u,v)$ is the distance.
In Erdős–Rényi random graphs above the connectivity threshold, the union of a minimum cut and a shortest path is within a factor $1+\epsilon$ of the optimum, for every fixed $\epsilon > 0$, with probability tending to one.

\end{abstract}

\begin{keywords}
NP-completeness \sep edge connectivity \sep graph diameter \sep Erdős–Rényi graphs \sep average-case approximation
\end{keywords}

\maketitle

% PDF author field: \maketitle of the class fills it with the author list followed by ", ", which the
% PDF shows as a stray character; set after \maketitle, before the first page is written.
\hypersetup{pdfauthor={Norbert Micheľ}}

\section{Introduction}
\label{sec:introduction}

Suppose that two trusted servers in a network must communicate while an adversary is denied every other route between them.
We model the network as a graph whose vertices are the servers and whose edges are the direct communication links.
The task is to select a set of edges that contains a connecting path, which carries the communication, and a separating cut, which guarantees that every other route between the two servers uses at least one of the selected edges.
We call such a set of edges a \emph{cut-path}.

Connectivity and separation are classical topics: multi-terminal minimum cuts~\cite{GomoryHu1961}, representations of all minimum cuts~\cite{NagamochiKameda1996}, and small cuts and edge connectivity~\cite{Mehlhorn2017Certifying} have been studied in detail, and paths and cuts are related by the max-flow min-cut duality~\cite{Cormen2022}.
These papers treat a path and a cut as separate objects.
A cut-path requires both in the same set of edges.


The \MinCutPath{} problem asks for a cut-path with the fewest edges between two given vertices.
To the best of our knowledge, the problem has not been studied before.
We settle its complexity and prove its basic structural properties.

Our contributions are the following.
First, we prove that the decision version of \MinCutPath{} is NP-complete, by a reduction from \ThreeSAT{} through an intermediate problem, \SSP.
Second, we prove that $\cp(u,v) = c(u,v) + d(u,v) - 1$ in graphs of diameter two and in graphs with minimum cut-value at most two.
Here $\cp(u,v)$ is the size of a minimum cut-path, $c(u,v)$ the size of a minimum cut, and $d(u,v)$ the distance; in both classes a minimum cut-path can be computed in polynomial time.
Third, dense Erdős–Rényi random graphs have diameter two with probability tending to one, so the formula applies to them.
For sparser random graphs, we show that the union of a minimum cut and a shortest path is within a factor of $1+\epsilon$ of the optimum on almost all inputs.

Section~\ref{sec:fundamentals} defines cut-paths and proves $\max\{c(u,v), d(u,v)\} \leq \cp(u,v) \leq c(u,v) + d(u,v) - 1$.
Section~\ref{sec:np-completeness} contains the NP-completeness proof.
Sections~\ref{sec:diameter-two} and~\ref{sec:cut-two} treat graphs of diameter two and graphs with minimum cut-value at most two, and Section~\ref{sec:random-graphs} treats random graphs.
Section~\ref{sec:conclusion} lists open problems.

\section{Fundamentals}
\label{sec:fundamentals}

Throughout the paper, we use standard notation from graph theory and computational complexity, following the conventions of~\cite{Diestel2025,Cormen2022}.
All graphs under consideration are undirected and finite, unless otherwise stated.
For a graph $G = (V, E)$, we let $n = |V|$ denote the number of vertices and $m = |E|$ the number of edges.

We often deal with two distinct vertices $u, v \in V$.
The value $d(u, v)$ denotes the distance between $u$ and $v$, i.e., the length of a shortest $u$--$v$ path in $G$.
Similarly, the value $c(u, v)$ denotes the minimum size of a cut that separates $u$ and $v$.
All problems studied in this paper are defined with respect to these two distinguished vertices $u$ and $v$, which are assumed to lie in the same connected component of $G$.

We identify a path with its set of edges.
For any path $P \subseteq E$, we denote by $G \setminus P$ the graph obtained from $G$ by removing all edges of $P$.
For a vertex $x \in V$ and a vertex subset $S \subseteq V$, we write $\deg(x, S)$ for the number of neighbors of $x$ that belong to $S$, that is, $\deg(x, S) = |\{\,y \in S : \{x, y\} \in E\,\}|$.



\begin{definition}[Cut-Path]
\label{def:cut-path}
Let $G = (V, E)$ be a graph, and let $u, v \in V$ be two distinct vertices.
A \emph{cut-path} between $u$ and $v$ is a set of edges $S \subseteq E$ for which there exist subsets $P, C \subseteq S$ such that:
    \begin{enumerate}
        \item $C$ is a cut separating $u$ and $v$ (i.e., removing $C$ disconnects $u$ from $v$),
        \item $P$ is a path connecting $u$ and $v$.
    \end{enumerate}
\end{definition}

Figure~\ref{fig:cut-path} shows a cut-path schematically.

\begin{figure}[H]
    \centering
    \includegraphics{cut-path.pdf}
    \caption{A cut-path between $u$ and $v$ (schematic)}
    \label{fig:cut-path}
\end{figure}

\begin{definition}[Cut-Paths]
\label{def:cut-paths}
Let $G = (V, E)$ be a graph, and let $u, v \in V$ be two distinct vertices.
The set of all cut-paths between $u$ and $v$, denoted by $\CP(u, v)$, is defined as:
\[
\CP(u, v) = \{ S \mid S \text{ is a cut-path between } u \text{ and } v \}.
\]
\end{definition}

\begin{definition}[Value of Minimum Cut-Path]
\label{def:cp-value}
Let $G = (V, E)$ be a graph, and let $u, v \in V$ be two distinct vertices.
The minimum size of a cut-path between $u$ and $v$, denoted by $\cp(u, v)$, is defined as:
\[
\cp(u, v) = \min\{ |S| \mid S \in \CP(u, v) \}.
\]
\end{definition}

A cut-path $S \in \CP(u, v)$ with $|S| = \cp(u, v)$ is called a \emph{minimum cut-path} between $u$ and $v$.
The corresponding optimization problem is stated as follows.

\begin{problem}[\MinCutPath{} (Optimization Version)]
Let $G=(V,E)$ be an undirected graph and $u,v \in V$ distinct vertices; cut-paths are defined in Definition~\ref{def:cut-path}.
The optimization version of the \MinCutPath{} problem is defined as follows:

\begin{center}
\begin{tabular}{p{0.18\linewidth}p{0.75\linewidth}}
\textbf{Input:} & A graph $G=(V,E)$ and vertices $u,v \in V$. \\
\textbf{Output:} & A minimum cut-path $S^* \subseteq E$ between $u$ and $v$, together with its cardinality
\[
\cp(u,v) = |S^*|.
\]
\end{tabular}
\end{center}
\end{problem}

The following elementary bounds relate $\cp(u, v)$ to the classical quantities $c(u, v)$ and $d(u, v)$; they are used repeatedly in the subsequent sections.

\begin{lemma}[Basic Bounds]
\label{lem:basic-bounds}
Let $G = (V, E)$ be a graph, and let $u, v \in V$ be two distinct vertices. Then
\[
\max\{c(u, v), d(u, v)\} \leq \cp(u, v) \leq c(u, v) + d(u, v) - 1.
\]
In particular, if $c(u, v) = 1$ or $d(u, v) = 1$, then $\cp(u, v) = c(u, v) + d(u, v) - 1$.
\end{lemma}

\begin{proof}
Every cut-path contains a $u$--$v$ cut, which has at least $c(u, v)$ edges, and a $u$--$v$ path, which has at least $d(u, v)$ edges; this gives the lower bound.
For the upper bound, let $C$ be a minimum $u$--$v$ cut and let $P$ be a shortest $u$--$v$ path.
The set $C \cup P$ is a cut-path.
Since removing $C$ disconnects $u$ from $v$, the path $P$ contains at least one edge of $C$, and hence $|C \cup P| \leq c(u, v) + d(u, v) - 1$.
If $c(u, v) = 1$ or $d(u, v) = 1$, the lower and the upper bound coincide.
\end{proof}


For random graphs, we use the following notion of approximation: a polynomial-time algorithm that is within a factor of $1+\epsilon$ of the optimum on almost all inputs.
Let $I(n)$ denote the set of all input instances of size $n$, and let $\OPT(x)$ denote the optimal value of the objective function for an instance $x \in I(n)$, i.e., the minimum for a minimization problem and the maximum for a maximization problem.

\begin{definition}[Average $(1+\epsilon)$-Approximation Scheme]
\label{def:approximation-scheme}
Let $\epsilon > 0$. An algorithm $A$ is called an \emph{average $(1+\epsilon)$-approximation scheme} if it meets the following criteria:
\begin{enumerate}
    \item $A$ is a polynomial-time algorithm, i.e., its running time is $O(P(n))$ for some polynomial $P(n)$.
    \item For each input size $n$, there exists a subset of inputs $I_{A}^{\mathrm{opt}}(n) \subseteq I(n)$ such that for every input $x \in I_{A}^{\mathrm{opt}}(n)$, the value $A(x)$ of the solution computed by $A$ satisfies
    \[
    \frac{A(x)}{\OPT(x)} < 1+\epsilon \quad \text{(for minimization problems)}
    \]
    \[
    \frac{\OPT(x)}{A(x)} < 1+\epsilon \quad \text{(for maximization problems)},
    \]
    where $\OPT(x)$ denotes the optimal value for the input $x$.
    \item In addition, the algorithm is average-case in the sense that
    \[
    \lim_{n \to \infty} \frac{\lvert I_{A}^{\mathrm{opt}}(n) \rvert}{\lvert I(n) \rvert} = 1.
    \]
\end{enumerate}
\end{definition}

